% Chapter 7, typeset from lectures/L07: the README is the opener and the closing review, bilaga A
% the sections, and bilaga B the exercises. The answers are in sv/back/facit/07.tex.

\chapter{Vektorer}
\label{c7:ch}

{\large\itshape Storheter med både storlek och riktning, och hur man räknar med dem.\par}

\begin{chapterbox}{I det här kapitlet}
\begin{itemize}
  \item Sträckor och vinklar i ett koordinatsystem.
  \item Vektorer: begrepp och notation.
  \item Absolutbelopp och vinkel, och vinkeln mellan två vektorer.
  \item Beräkningar med vektorer: addition, subtraktion och multiplikation med ett tal.
\end{itemize}
\end{chapterbox}

Ett tal som en resistans eller en temperatur har bara en storlek. Många storheter har också en
riktning: en kraft drar åt ett visst håll, och i växelströmskretsar beskrivs spänningar och
strömmar av en storlek och en fasvinkel. Sådana storheter beskrivs med vektorer. Det här kapitlet
börjar med sträckor och vinklar i ett koordinatsystem och går sedan igenom vektorer: hur de
skrivs med komponenter, hur deras längd och vinkel beräknas och hur de adderas, subtraheras och
multipliceras med ett tal. Vinkeln mellan två vektorer fås ur vektorernas egna vinklar. Vektorerna
återkommer senare i boken, när spänningar och strömmar räknas som visare och komplexa tal.

\Secref{c7:sec:koordinatsystem} till \secref{c7:sec:rakna} är kapitlets teori, och
\secref{c7:sec:typexempel} visar den i ett typexempel. \Secref{c7:sec:sammanfattning}
sammanfattar formlerna och vad du ska kunna efter kapitlet, och i \secref{c7:sec:uppgifter} finns
uppgifterna.

% Sections 7.1 to 7.6, from lectures/L07/appendix/a_theory.md.

\section{Koordinatsystem och sträckor}
\label{c7:sec:koordinatsystem}
Ett \textbf{koordinatsystem} ($xy$-plan) har en horisontell $x$-axel och en vertikal $y$-axel
som möts i origo $(0, 0)$.

Sträckan $d$ mellan två punkter $(x_1, y_1)$ och $(x_2, y_2)$ ges av Pythagoras sats:
\[
  d = \sqrt{(x_2 - x_1)^2 + (y_2 - y_1)^2}
\]
Vinkeln som sträckan bildar med den positiva $x$-axeln är
\[
  \theta = \arctan\left(\frac{y_2 - y_1}{x_2 - x_1}\right).
\]

\begin{aside}
\note[Obs] Kontrollera alltid vilken kvadrant punkten befinner sig i och korrigera vinkeln vid
behov: $+180^{\circ}$ för kvadrant II och III.
\end{aside}

\section{Vektorer -- begrepp och notation}
\label{c7:sec:begrepp}
En \textbf{vektor} representerar en storhet med både storlek och riktning. I ett plan skrivs
vektorn $\mathbf{u}$ med sina komponenter:
\[
  \mathbf{u} = (u_x;\, u_y)
\]
Till skillnad från ett skalärt tal bär vektorn information om riktning. Vektorer används inom
elektroteknik för att representera spänningar och strömmar i AC-system.

\section{Absolutbelopp och vinkel}
\label{c7:sec:absolutbelopp}
\textbf{Absolutbeloppet} (längden) av $\mathbf{u} = (u_x;\, u_y)$ är
\[
  \abs{\mathbf{u}} = \sqrt{u_x^2 + u_y^2}.
\]
\textbf{Vinkeln} $v_u$ mot den positiva $x$-axeln är
\[
  v_u = \arctan\left(\frac{u_y}{u_x}\right) + \text{kvadrantkorrigering},
\]
där kvadrantkorrigeringen ges av tabellen.

\begin{narrowtable}{@{}cccl@{}}
  \thead{Kvadrant} & $u_x$ & $u_y$ & \thead{Korrigering} \\
  \midrule
  I & $+$ & $+$ & ingen \\
  II & $-$ & $+$ & $+180^{\circ}$ \\
  III & $-$ & $-$ & $+180^{\circ}$ \\
  IV & $+$ & $-$ & ingen (negativ vinkel) \\
\end{narrowtable}

\section{Vinkeln mellan två vektorer}
\label{c7:sec:vinkel}
Vinkeln $\alpha$ mellan $\mathbf{u}$ och $\mathbf{v}$ är skillnaden mellan vektorernas vinklar mot
den positiva $x$-axeln. Beräkna därför först $v_u$ och $v_v$ var för sig och därefter differensen:
\[
  \alpha = v_v - v_u
\]
Vinkeln mellan två vektorer anges som ett värde mellan $0^{\circ}$ och $180^{\circ}$. Blir
differensen negativ byter du tecken på den, och blir den större än $180^{\circ}$ drar du den från
$360^{\circ}$.

\section{Räkna med vektorer}
\label{c7:sec:rakna}

\subsection{Addition och subtraktion}
\label{c7:sec:addition}
\[
  \mathbf{u} \pm \mathbf{v} = (u_x \pm v_x;\; u_y \pm v_y)
\]

\subsection{Multiplikation med ett tal}
\label{c7:sec:multiplikation}
När en vektor multipliceras med ett tal $k$ multipliceras båda komponenterna med talet:
\[
  k\mathbf{u} = (k \cdot u_x;\; k \cdot u_y)
\]
Exempelvis gäller för $\mathbf{U} = (3;\, 4)$ att
\[
  2\mathbf{U} = (2 \cdot 3;\; 2 \cdot 4) = (6;\, 8).
\]
Ett negativt tal byter dessutom riktning på vektorn.

\section{Typexempel}
\label{c7:sec:typexempel}

\begin{exempel}[Vektorer i ett koordinatsystem]
Två vektorer $\mathbf{u}$ samt $\mathbf{v}$ visas i nedanstående figur.

\bild[0.4\linewidth]{lectures/L07/appendix/images/example1_1.png}

\begin{parts}
  \item Ta reda på vektorernas koordinater.
  \item Beräkna vektorernas absolutbelopp $\abs{\mathbf{u}}$ samt $\abs{\mathbf{v}}$.
  \item Beräkna vektorernas vinklar $v_u$ samt $v_v$.
  \item Beräkna vinkeln mellan vektorer $\mathbf{u}$ samt $\mathbf{v}$.
  \item Beräkna en tredje vektor $\mathbf{w} = \mathbf{u} + 2\mathbf{v}$.
  \item Bestäm en vektor med längden $5$ som är motsatt riktad $\mathbf{u}$.
\end{parts}

\textbf{a)} Genom att kolla i figuren ser vi att vektor $\mathbf{u}$ har $x$-koordinaten $3$
samt $y$-koordinaten $3$. Därmed gäller att
\[
  \mathbf{u} = (3;\, 3)
\]
Vi ser också att vektor $\mathbf{v}$ har $x$-koordinaten $-2$ samt $y$-koordinaten $3$. Därmed
gäller att
\[
  \mathbf{v} = (-2;\, 3)
\]

\textbf{b)} Vektorernas absolutbelopp är lika med deras respektive hypotenusa och beräknas
därmed enkelt med Pythagoras sats.

Detta kan visualiseras genom att rita ut respektive vektor som en triangel, se figuren nedan.

\bild[0.4\linewidth]{lectures/L07/appendix/images/example1_2.png}

Vektorernas absolutbelopp $\abs{\mathbf{u}}$ samt $\abs{\mathbf{v}}$ beräknas enligt nedan:
\begin{gather*}
  \abs{\mathbf{u}} = \sqrt{3^2 + 3^2} = \sqrt{18} \approx 4{,}24 \\
  \abs{\mathbf{v}} = \sqrt{3^2 + (-2)^2} = \sqrt{13} \approx 3{,}61
\end{gather*}

\textbf{c)} För att beräkna vektorernas vinklar kan vi använda oss utav tangens. Återigen ser vi
vektorerna som trianglar.

För vektor $\mathbf{u} = (3;\, 3)$ gäller att både motstående samt närliggande katet är lika
med $3$. Därmed gäller att
\[
  \tan(v_u) = \frac{3}{3} = 1,
\]
vilket kan transformeras till
\[
  v_u = \arctan(1) = 45^{\circ}
\]
För vektor $\mathbf{v}$, som har ett negativt $x$-värde, kan motsvarande triangel ritas ut såsom
visas i ovanstående figur. Notera att vi i detta fall har markerat en vinkel $v_x$, som är den
supplementära vinkeln till $v_v$:
\[
  v_v + v_x = 180^{\circ},
\]
vilket innebär att
\[
  v_v = 180^{\circ} - v_x
\]
Om vi beräknar vinkel $v_x$ kan vi därmed enkelt beräkna vinkel $v_v$. Vinkel $v_x$ beräknas
enkelt med tangens, då motsvarande triangel har en motstående katet som är lika med $3$ samt en
närliggande katet som är lika med $2$. Därmed gäller att
\[
  \tan(v_x) = \frac{3}{2} = 1{,}5 \quad \Rightarrow \quad
  v_x = \arctan(1{,}5) \approx 56{,}3^{\circ}
\]
Vi kan sedan enkelt beräkna vinkel $v_v$:
\[
  v_v = 180^{\circ} - v_x = 180^{\circ} - 56{,}3^{\circ} = 123{,}7^{\circ}
\]

\textbf{d)} För att beräkna vinkeln $\alpha$ mellan $\mathbf{v}$ samt $\mathbf{u}$ beräknar vi
differensen mellan deras respektive vinklar $v_v$ samt $v_u$:
\[
  \alpha = v_v - v_u
\]
Vi sätter in värden i formeln och ser då att vinkeln mellan vektorerna är ca $78{,}7^{\circ}$, då
\[
  \alpha = v_v - v_u = 123{,}7^{\circ} - 45^{\circ} \approx 78{,}7^{\circ}
\]

\textbf{e)} Vi beräknar vektor $\mathbf{w}$ algebraiskt genom att sätta in värdena av vektorer
$\mathbf{u} = (3;\, 3)$ samt $\mathbf{v} = (-2;\, 3)$ i given formel:
\[
  \mathbf{w} = \mathbf{u} + 2\mathbf{v} = (3;\, 3) + 2(-2;\, 3),
\]
där
\[
  2(-2;\, 3) = (2 \cdot (-2);\; 2 \cdot 3) = (-4;\, 6)
\]
Därmed gäller att
\[
  \mathbf{w} = (3;\, 3) + (-4;\, 6)
\]
För att slutföra beräkningen summerar vi ihop $x$- och $y$-värdena och ser då att
$\mathbf{w} = (-1;\, 9)$, då
\[
  \mathbf{w} = (3 - 4;\; 3 + 6) = (-1;\, 9)
\]

\textbf{f)} Vi kan rita ut en vektor motsatt riktad $\mathbf{u}$ med längden $5$. Låt oss kalla
denna vektor $\mathbf{a}$. Att vektorn är motriktad innebär att den är vänd $180^{\circ}$,
alltså ett halvt varv från $\mathbf{u}$. För dess vinkel $v_a$ gäller därmed att
\[
  v_a = v_u + 180^{\circ}
\]
Eftersom vinkel $v_u$ tidigare beräknades till $45^{\circ}$ gäller att vektor $\mathbf{a}$:s
vinkel $v_a$ är lika med $225^{\circ}$, då
\[
  v_a = 45^{\circ} + 180^{\circ} = 225^{\circ}
\]
För en given vektor gäller att dess längd är lika med dess absolutbelopp. Därmed gäller att
absolutbeloppet $\abs{\mathbf{a}}$ av vektorn $\mathbf{a}$ är lika med $5$:
\[
  \abs{\mathbf{a}} = 5
\]
Vi kan sedan beräkna vektor $\mathbf{a}$ via två metoder, grafiskt eller algebraiskt.

\textbf{Metod 1 -- Grafisk lösning.} Som vi såg tidigare gäller att vektorn kan ritas ut som en
triangel, där dess hypotenusa är lika med dess absolutbelopp. Eftersom denna vektor är vänd ett
halvt varv från vektor $\mathbf{u}$ kan vi rita vektor $\mathbf{a}$ som en triangel med vinkel
$45^{\circ}$ sett från $180^{\circ}$, dvs. vänt ett halvt varv.

Med denna information kan vi rita ut vektor $\mathbf{a}$ såsom visas nedan.

\bild[0.4\linewidth]{lectures/L07/appendix/images/example1_3.png}

Vi kan beräkna absolutbeloppet $\abs{x_a}$ av vektor $\mathbf{a}$:s närliggande katet $x_a$ med
cosinus:
\[
  \cos 45^{\circ} = \frac{\abs{x_a}}{5},
\]
vilket innebär att
\[
  \abs{x_a} = 5 \cdot \cos 45^{\circ} \approx 3{,}54
\]
Eftersom $x_a$ är placerad till vänster om origo är $x_a$ negativ. Därmed gäller att
\[
  x_a \approx -3{,}54
\]
Vi kan också beräkna absolutbeloppet $\abs{y_a}$ av vektor $\mathbf{a}$:s motstående katet
$y_a$ med sinus:
\[
  \sin 45^{\circ} = \frac{\abs{y_a}}{5},
\]
vilket innebär att
\[
  \abs{y_a} = 5 \cdot \sin 45^{\circ} \approx 3{,}54
\]
Eftersom $y_a$ är placerad under origo är $y_a$ negativ. Därmed gäller att
\[
  y_a \approx -3{,}54
\]
Via beräknade värden kan sedan vektor $\mathbf{a}$ bestämmas:
\[
  \mathbf{a} \approx (-3{,}54;\; -3{,}54)
\]

\textbf{Metod 2 -- Algebraisk lösning.} Vi såg tidigare att vektor $\mathbf{a}$:s vinkel $v_a$
är lika med $225^{\circ}$, då denna är vänd ett halvt varv ($180^{\circ}$) från vektor
$\mathbf{u}$, vars vinkel $v_u$ tidigare har beräknats till $45^{\circ}$:
\[
  v_a = v_u + 180^{\circ} = 45^{\circ} + 180^{\circ} = 225^{\circ}
\]
Vektorns absolutbelopp $\abs{\mathbf{a}}$, som är lika med dess längd, specificerades också till
$5$ i uppgiftsbeskrivningen:
\[
  \abs{\mathbf{a}} = 5
\]
Vi kan enkelt beräkna vektorns $x$- och $y$-värde via cosinus samt sinus. Vi börjar med att
beräkna vektorns $x$-värde $x_a$:
\[
  \cos v_a = \frac{x_a}{\abs{\mathbf{a}}},
\]
vilket är ekvivalent med att
\[
  x_a = \abs{\mathbf{a}} \cdot \cos v_a
\]
Genom att sätta in värden i ovanstående formel ser vi att vektorns $x$-värde $x_a$ är lika med
$-3{,}54$, då
\[
  x_a = 5 \cdot \cos 225^{\circ} \approx -3{,}54
\]
Vi kan sedan beräkna vektor $\mathbf{a}$:s $y$-värde $y_a$ med sinus:
\[
  \sin v_a = \frac{y_a}{\abs{\mathbf{a}}},
\]
vilket är ekvivalent med att
\[
  y_a = \abs{\mathbf{a}} \cdot \sin v_a
\]
Genom att sätta in värden i ovanstående formel ser vi att även vektor $\mathbf{a}$:s $y$-värde
$y_a$ är lika med $-3{,}54$, då
\[
  y_a = 5 \cdot \sin 225^{\circ} \approx -3{,}54
\]
Därmed gäller att
\[
  \mathbf{a} \approx (-3{,}54;\; -3{,}54)
\]
\end{exempel}

% Section 7.7, from the summary table of lectures/L07/appendix/a_theory.md and from
% lectures/L07/README.md: the objectives, the evaluation questions and the next lesson.

\section{Sammanfattning}
\label{c7:sec:sammanfattning}

\begin{booktable}{@{}p{12.5em}L@{}}
  \thead{Begrepp} & \thead{Formel} \\
  \midrule
  Absolutbelopp & $\abs{\mathbf{u}} = \sqrt{u_x^2 + u_y^2}$ \\
  Vinkel mot $x$-axeln & $v = \arctan(u_y / u_x)$ + kvadrantkorrigering \\
  Vinkel mellan vektorer & $\alpha = v_v - v_u$ \\
  Addition & $\mathbf{u} + \mathbf{v} = (u_x + v_x;\, u_y + v_y)$ \\
  Multiplikation med ett tal & $k\mathbf{u} = (ku_x;\, ku_y)$ \\
\end{booktable}

\needspace{6\baselineskip}
\subsection*{Det här ska du kunna}
Efter det här kapitlet ska du:
\begin{itemize}
  \item Förstå begreppet vektor och hur en vektor representeras matematiskt.
  \item Kunna genomföra enkla beräkningar med vektorer.
\end{itemize}

\testadigsjalv
\begin{enumerate}
  \item En vektor är $\mathbf{u} = (3, -4)$. Beräkna absolutbeloppet $\abs{\mathbf{u}}$ och
    vinkeln mot $x$-axeln.
  \item Beräkna $\mathbf{w} = 2\mathbf{u} - \mathbf{v}$ för $\mathbf{u} = (2, 1)$ och
    $\mathbf{v} = (-1, 3)$.
\end{enumerate}

\needspace{6\baselineskip}
\subsection*{Nästa kapitel}
\Kapref{8} handlar om funktioner.

% Section 7.8, from lectures/L07/appendix/b_exercises.md.

\section{Uppgifter}
\label{c7:sec:uppgifter}

\begin{aside}
\textbf{Så arbetar du med uppgifterna.} Del~1 och del~2 räknas under lektionen: först på egen
hand, några minuter per uppgift, sedan gemensamt. Del~3 är extrauppgifter att träna vidare på
efter lektionen. De flesta går att räkna i huvudet eller med papper och penna.

\facitnotis{L07}
\end{aside}

% ------------------------------------------------------------------------------------------
\uppgiftsdel{Del 1}{Repetitionsuppgifter}

\uppgift{1.1}{Sensormodul på en järnvägsvagn}
En sensormodul ute på en järnvägsvagn matas från vagnens $12\,\text{V}$-system. För att begränsa
strömmen vid kortslutning sitter ett skyddsmotstånd på $4\,\text{k}\Omega$ i serie med modulen.
Vid en mätning visar det sig att sensormodulen förbrukar $5\,\text{mW}$.

\bild[0.5\linewidth]{lectures/L07/appendix/images/1.1_circuit.png}

Du vill ta reda på hur stor strömmen $I$ genom kretsen är, men har ingen amperemeter till hands,
bara dina kunskaper i ellära.
\begin{parts}
  \item Spänningen över skyddsmotståndet är $4I$ (med $I$ i mA ger
    $\text{k}\Omega \cdot \text{mA} = \text{V}$). Visa att spänningen över sensormodulen blir
    $12 - 4I$ och att effektsambandet leder till ekvationen
    \[
      4I^2 - 12I + 5 = 0.
    \]
  \item Lös ekvationen för strömmen $I$.
  \item Ekvationen ger två lösningar. Beräkna spänningen över sensormodulen och effektförlusten i
    skyddsmotståndet för båda fallen. Vilken driftpunkt är mest rimlig för en sensormodul, och
    varför?
\end{parts}

% ------------------------------------------------------------------------------------------
\uppgiftsdel{Del 2}{Nytt stoff}
I växelströmskretsar (AC) representeras spänningar och strömmar av vektorer, så kallade
\textbf{visare}: vektorns längd anger storleken och dess vinkel mot den positiva $x$-axeln anger
fasvinkeln. Två visare med samma riktning är \emph{i fas}, och två motriktade visare är \emph{i
motfas}.

\uppgift{2.1}{Addition av spänningar}
Två växelspänningar representeras av visarna $\mathbf{U}_1 = (3;\, 2)\,\text{V}$ och $\mathbf{U}_2
= (4;\, 6)\,\text{V}$. Bestäm koordinaterna för:
\begin{delar}(4)
  \task $\mathbf{U}_1 + \mathbf{U}_2$
  \task $2\mathbf{U}_1 + 3\mathbf{U}_2$
  \task $\mathbf{U}_1 - \mathbf{U}_2$
  \task $-2\mathbf{U}_1 - 4\mathbf{U}_2$
\end{delar}

\uppgift{2.2}{Räkning med tre strömmar}
Tre växelströmmar representeras av visarna $\mathbf{I}_a = (1;\, 4)\,\text{A}$, $\mathbf{I}_b =
(-4;\, 2)\,\text{A}$ och $\mathbf{I}_c = (3;\, -4)\,\text{A}$. Bestäm koordinaterna för:
\begin{delar}(4)
  \task $\mathbf{I}_a + \mathbf{I}_b + \mathbf{I}_c$
  \task $\mathbf{I}_a - \mathbf{I}_b - \mathbf{I}_c$
  \task $2\mathbf{I}_a + 4\mathbf{I}_b$
  \task $2\mathbf{I}_a - 4\mathbf{I}_b$
\end{delar}

\uppgift{2.3}{Absolutbelopp av spänningar}
Spänningarna $\mathbf{U}_a = (2;\, 1)\,\text{V}$ och $\mathbf{U}_b = (1;\, -2)\,\text{V}$ är
givna. Bestäm i exakt form storleken (absolutbeloppet) av:
\begin{delar}(4)
  \task $\mathbf{U}_a$
  \task $3\mathbf{U}_a + \mathbf{U}_b$
  \task $\mathbf{U}_a - \mathbf{U}_b$
  \task $3\mathbf{U}_a - 4\mathbf{U}_b$
\end{delar}

\uppgift{2.4}{Lös ekvationerna}
Bestäm $x$ och $y$:
\begin{delar}(2)
  \task $(x;\, y) = 5(3;\, 1) + 2(4;\, -5)$
  \task $(15;\, 12) = (x;\, 10) + (23;\, y)$
\end{delar}

\uppgift{2.5}{Strömmar i fas och i motfas}
Om $\mathbf{I}_1 = k \cdot \mathbf{I}_2$ är visarna $\mathbf{I}_1$ och $\mathbf{I}_2$ parallella:
i fas om $k > 0$ och i motfas om $k < 0$. Vilka av följande strömmar är parallella med $\mathbf{I}
= (2;\, -3)\,\text{A}$?
\begin{delar}(4)
  \task $(6;\, -9)\,\text{A}$
  \task $(-8;\, 16)\,\text{A}$
  \task $(10;\, 15)\,\text{A}$
  \task $(-12;\, 18)\,\text{A}$
\end{delar}

\uppgift{2.6}{Bestäm okända koefficienter}
Spänningarna $\mathbf{U}_1 = (-2;\, 3)\,\text{V}$ och $\mathbf{U}_2 = (3;\, -1)\,\text{V}$ är
givna. Bestäm talen $a$ och $b$ så att
\[
  a\mathbf{U}_1 + \mathbf{U}_2 = (0;\, b)\,\text{V}.
\]

\uppgift{2.7}{Fasförskjutning mellan spänning och ström}
Spänningen över en last är $\mathbf{U} = (4;\, 4)\,\text{V}$ och strömmen genom den är $\mathbf{I}
= (-1;\, 2)\,\text{A}$.
\begin{parts}
  \item Beräkna storlekarna $\abs{\mathbf{U}}$ och $\abs{\mathbf{I}}$.
  \item Beräkna fasvinklarna $v_U$ och $v_I$ mot den positiva $x$-axeln.
  \item Beräkna fasförskjutningen, alltså vinkeln $\alpha$ mellan $\mathbf{U}$ och $\mathbf{I}$.
\end{parts}

% ------------------------------------------------------------------------------------------
\uppgiftsdel{Del 3}{Extrauppgifter}
Uppgifterna nedan är extra träning och görs med fördel på egen hand efter lektionen.

\uppgift{3.1}{Avstånd mellan punkter}
Beräkna avståndet mellan punkterna:
\begin{delar}(3)
  \task $(0, 0)$ och $(3, 4)$
  \task $(1, 2)$ och $(4, 6)$
  \task $(-2, 1)$ och $(2, -2)$
  \task $(5, -1)$ och $(5, 7)$
  \task $(-3, 2)$ och $(1, 2)$
\end{delar}

\uppgift{3.2}{Vektorn mellan två punkter}
Vektorn från punkten $A$ till punkten $B$ ges av $\overrightarrow{AB} = (x_B - x_A;\, y_B - y_A)$.
\begin{parts}
  \item $A = (1, 2)$ och $B = (5, 5)$. Bestäm $\overrightarrow{AB}$ och $\abs{\overrightarrow{AB}}$.
  \item $A = (-1, 4)$ och $B = (2, 0)$. Bestäm $\overrightarrow{AB}$ och $\abs{\overrightarrow{AB}}$.
  \item Bestäm $\overrightarrow{BA}$ i a) och jämför med $\overrightarrow{AB}$.
  \item Vad gäller allmänt mellan $\overrightarrow{AB}$ och $\overrightarrow{BA}$?
\end{parts}

\uppgift{3.3}{Spänningens storlek}
Bestäm spänningens storlek $\abs{\mathbf{U}}$ i exakt form och som närmevärde med två decimaler.
\begin{delar}(3)
  \task $\mathbf{U} = (6;\, 8)\,\text{V}$
  \task $\mathbf{U} = (-5;\, 12)\,\text{V}$
  \task $\mathbf{U} = (1;\, 1)\,\text{V}$
  \task $\mathbf{U} = (-3;\, -3)\,\text{V}$
  \task $\mathbf{U} = (0;\, -7)\,\text{V}$
\end{delar}

\uppgift{3.4}{Fasvinkel}
Bestäm strömmens fasvinkel $v$ mot den positiva $x$-axeln. Ange vilken kvadrant visaren ligger i
och om kvadrantkorrigering behövs.
\begin{delar}(3)
  \task $\mathbf{I} = (1;\, 1)\,\text{A}$
  \task $\mathbf{I} = (-1;\, 1)\,\text{A}$
  \task $\mathbf{I} = (-1;\, -1)\,\text{A}$
  \task $\mathbf{I} = (1;\, -1)\,\text{A}$
  \task $\mathbf{I} = (4;\, 3)\,\text{A}$
\end{delar}

\uppgift{3.5}{Från storlek och fasvinkel till komponenter}
En spänning med storleken $\abs{\mathbf{U}}$ och fasvinkeln $v$ mot den positiva $x$-axeln har
komponenterna
\[
  U_x = \abs{\mathbf{U}} \cos v, \qquad U_y = \abs{\mathbf{U}} \sin v.
\]
Bestäm komponenterna:
\begin{delar}(2)
  \task $\abs{\mathbf{U}} = 10\,\text{V}$, $v = 30^{\circ}$
  \task $\abs{\mathbf{U}} = 5\,\text{V}$, $v = 90^{\circ}$
  \task $\abs{\mathbf{U}} = 4\,\text{V}$, $v = 180^{\circ}$
  \task $\abs{\mathbf{U}} = 2\,\text{V}$, $v = 225^{\circ}$
  \task* Kontrollera i a) att $\sqrt{U_x^2 + U_y^2} = 10\,\text{V}$.
\end{delar}

\uppgift{3.6}{Fasförskjutning}
Beräkna först fasvinklarna $v_U$ och $v_I$ mot den positiva $x$-axeln och därefter
fasförskjutningen $\alpha$, alltså vinkeln mellan spänningen och strömmen.
\begin{delar}(2)
  \task $\mathbf{U} = (1;\, 0)\,\text{V}$, $\mathbf{I} = (1;\, 1)\,\text{A}$
  \task $\mathbf{U} = (3;\, 0)\,\text{V}$, $\mathbf{I} = (0;\, 4)\,\text{A}$
  \task $\mathbf{U} = (2;\, 1)\,\text{V}$, $\mathbf{I} = (-1;\, 3)\,\text{A}$
  \task $\mathbf{U} = (2;\, 2)\,\text{V}$, $\mathbf{I} = (-3;\, -3)\,\text{A}$
\end{delar}

\uppgift{3.7}{Motstånd och spole i serie}
Över ett motstånd och en spole i serie ligger spänningarna $\mathbf{U}_R = (12;\, 0)\,\text{V}$
och $\mathbf{U}_L = (0;\, 5)\,\text{V}$.
\begin{parts}
  \item Beräkna den totala spänningen $\mathbf{U} = \mathbf{U}_R + \mathbf{U}_L$.
  \item Beräkna storleken $\abs{\mathbf{U}}$.
  \item Beräkna fasvinkeln mot den positiva $x$-axeln.
  \item Varför är $\abs{\mathbf{U}}$ inte lika med $12 + 5 = 17\,\text{V}$?
\end{parts}

\uppgift{3.8}{Motstånd, spole och kondensator i serie}
I en seriekrets ligger spänningen $\mathbf{U}_R = (30;\, 0)\,\text{V}$ över motståndet,
$\mathbf{U}_L = (0;\, 70)\,\text{V}$ över spolen och $\mathbf{U}_C = (0;\, -30)\,\text{V}$ över
kondensatorn.
\begin{parts}
  \item Beräkna den totala spänningen $\mathbf{U} = \mathbf{U}_R + \mathbf{U}_L + \mathbf{U}_C$.
  \item Beräkna storleken $\abs{\mathbf{U}}$.
  \item Beräkna fasvinkeln mot den positiva $x$-axeln.
  \item Vad blir den totala spänningen om kondensatorns spänning i stället är $(0;\, -70)\,\text{V}$?
    Vad innebär resultatet?
\end{parts}

\uppgift{3.9}{Enhetsvektorer}
En enhetsvektor har längden $1$ och fås som $\hat{\mathbf{u}} =
\dfrac{\mathbf{u}}{\abs{\mathbf{u}}}$. Strömmen $\mathbf{I} = (3;\, 4)\,\text{A}$ är given.
\begin{parts}
  \item Bestäm enhetsvektorn i strömmens riktning.
  \item Kontrollera att svaret i a) verkligen har längden $1$.
  \item Bestäm enhetsvektorn i riktning $(-6;\, 8)$.
  \item Bestäm en ström med storleken $15\,\text{A}$ som är i fas med $\mathbf{I}$.
  \item Bestäm en ström med storleken $15\,\text{A}$ som är i motfas med $\mathbf{I}$.
\end{parts}

\uppgift{3.10}{Trefassystem}
I ett trefassystem har de tre fasspänningarna samma storlek men olika fasvinklar: $\mathbf{U}_1 =
(230;\, 0)\,\text{V}$, $\mathbf{U}_2 = (-115;\, -199)\,\text{V}$ och $\mathbf{U}_3 = (-115;\,
199)\,\text{V}$.
\begin{parts}
  \item Beräkna summan $\mathbf{U}_1 + \mathbf{U}_2 + \mathbf{U}_3$.
  \item Vad innebär resultatet?
  \item Beräkna storleken och fasvinkeln för $\mathbf{U}_2$.
  \item Beräkna huvudspänningen $\mathbf{U}_{12} = \mathbf{U}_1 - \mathbf{U}_2$ och dess storlek.
\end{parts}

\uppgift{3.11}{I fas och i motfas}
Spänningen $\mathbf{U} = (4;\, -6)\,\text{V}$ är given.
\begin{parts}
  \item Vilka av $(2;\, -3)\,\text{V}$, $(-8;\, 12)\,\text{V}$, $(6;\, -4)\,\text{V}$ och $(-2;\,
    3)\,\text{V}$ är parallella med $\mathbf{U}$?
  \item Vilka av dem är i motfas med $\mathbf{U}$?
  \item Bestäm $t$ så att $(t;\, 9)\,\text{V}$ blir parallell med $\mathbf{U}$.
  \item Kan en visare vara parallell med sig själv? Motivera.
\end{parts}

\uppgift{3.12}{Motstånd och kondensator parallellt}
Genom ett motstånd och en kondensator som är parallellkopplade flyter strömmarna $\mathbf{I}_R =
(4;\, 0)\,\text{A}$ och $\mathbf{I}_C = (0;\, 3)\,\text{A}$.
\begin{parts}
  \item Beräkna den totala strömmen $\mathbf{I} = \mathbf{I}_R + \mathbf{I}_C$.
  \item Beräkna storleken $\abs{\mathbf{I}}$.
  \item Beräkna fasvinkeln mot den positiva $x$-axeln.
  \item Hur stor måste kondensatorströmmen vara för att fasvinkeln ska bli $45^{\circ}$?
\end{parts}

\uppgift{3.13}{Effekttriangeln}
Effekten i en växelströmskrets kan beskrivas med vektorn $\mathbf{S} = (P;\, Q)$, där $P$ är den
aktiva effekten i watt (W) och $Q$ är den reaktiva effekten i var. Vektorns storlek
$\abs{\mathbf{S}}$ är den skenbara effekten i voltampere (VA). En motor har $P = 800\,\text{W}$
och $Q = 600\,\text{var}$.
\begin{parts}
  \item Beräkna den skenbara effekten $\abs{\mathbf{S}}$.
  \item Beräkna vektorns vinkel $v$ mot den positiva $x$-axeln.
  \item Effektfaktorn är $\cos v$. Beräkna den.
  \item En kondensator kopplas in och bidrar med vektorn $(0;\, -600)$. Bestäm den nya vektorn
    $\mathbf{S}$ och den nya skenbara effekten.
\end{parts}

\uppgift{3.14}{Hitta felet}
Varje rad innehåller ett vanligt vektorfel. Förklara felet och ange det korrekta svaret.
\begin{parts}
  \item $\abs{(3;\, 4)} = 3 + 4 = 7$
  \item $2(3;\, -1) = (6;\, -1)$
  \item Två visare har vinklarna $30^{\circ}$ och $100^{\circ}$ mot den positiva $x$-axeln, så vinkeln
    mellan dem är $130^{\circ}$.
  \item Vektorn $(-3;\, 4)$ bildar vinkeln $\arctan\!\left(\dfrac{4}{-3}\right) \approx -53{,}1^{\circ}$
    mot den positiva $x$-axeln.
  \item $\abs{-2\mathbf{u}} = -2\abs{\mathbf{u}}$
\end{parts}

